\chapter{Adaptación de los NAS Parallel Benchmarks a OMP4Py}
\label{chap:adaptacion-nas}

\section{Benchmarks adaptados}
\label{sec:benchmarks-adaptados}

Se adaptaron cinco NAS Parallel Benchmarks al modelo de programación OMP4Py,
en concreto EP, CG, IS, MG y FT. En todos los casos se conservaron la clase del problema,
el número de iteraciones y la verificación oficial del benchmark como criterio
único de aceptación. Los tres benchmarks restantes del conjunto NPB (BT, SP y
LU) quedan diferidos por dos motivos distintos. BT y SP se apoyan en el método
ADI, con resoluciones tridiagonales a lo largo de cada eje. Como las líneas son
independientes entre sí, su paralelización encaja en un \texttt{parallel for}
ordinario, y quedaron fuera sobre todo por el tiempo disponible. El caso de LU
es distinto.
Su barrido SSOR (rutinas \texttt{blts} y \texttt{buts}) arrastra una dependencia
en \emph{wavefront}, porque cada punto necesita a sus vecinos ya actualizados en
las tres direcciones, así que no admite un bucle paralelo directo. Las
formulaciones habituales en OpenMP lo resuelven de tres maneras, con el
constructo \texttt{ordered} (bucles \emph{doacross}), con sincronización punto a
punto mediante \texttt{flush} y \texttt{atomic}, o con tareas y dependencias
(\texttt{depend}). El conjunto de directivas que expone OMP4Py en su versión
actual no incluye ninguno de esos mecanismos, así que un port directo del LU
paralelo de NPB no llega a ser expresable. Solo cabría reescribirlo por
hiperplanos apoyándose en \texttt{parallel for} y \texttt{barrier}, que supone
un trabajo algorítmico aparte. En las adaptaciones actuales, los tres
benchmarks se mantienen en versión serial, con compilación Cython pero sin
directivas OMP4Py.

Todos los benchmarks siguen una misma interfaz de línea de comandos, basada en
las opciones \texttt{-c} para la clase NAS, \texttt{-t} para el número de hilos y
\texttt{-m} para el modo de ejecución. La invocación completa se recoge en el
Anexo~\ref{anexo:interfaz}.

La Tabla~\ref{tab:estado-benchmarks} resume el estado funcional de los ocho
benchmarks disponibles en el repositorio.

\begin{table}[htbp]
\centering
\begin{tabular}{l|l|l}
Benchmark & Estado & Observación principal \\
\hline
EP & Completo & Reducción paralela y variables privadas tipadas \\
CG & Completo & Producto disperso y reducciones escalares \\
IS & Completo & Semántica de cubos revisada para OMP4Py \\
MG & Completo & Kernels de malla con memoryviews 1D \\
FT & Completo & FFT compleja con región paralela exterior \\
BT & Diferido & Sin directivas OMP4Py, solo compilación Cython \\
SP & Diferido & Sin directivas OMP4Py, solo compilación Cython \\
LU & Diferido & Sin directivas OMP4Py, solo compilación Cython \\
\end{tabular}
\caption{Estado de los benchmarks NAS en el repositorio.}
\label{tab:estado-benchmarks}
\end{table}


\section{Metodología de adaptación}
\label{sec:metodologia-portado}

El port se realizó de forma incremental, mediante un ciclo de portar, ejecutar,
verificar y depurar. Se partió de las versiones Python seriales de
referencia~\cite{NPBPython}, derivadas de los NAS en C++, y cuando existía se
consultó la versión C/OpenMP
original para entender el patrón de paralelismo. Los pasos seguidos fueron los
siguientes.

\begin{enumerate}
  \item Partir de la versión Python serial NAS.
  \item Comparar con la versión C/OpenMP cuando existe.
  \item Conservar clase, tamaño, iteraciones, salida y verificación NAS.
  \item Eliminar dependencias de Numba~\cite{Numba2015} y adaptar imports.
  \item Añadir la interfaz común \texttt{-c}, \texttt{-t}, \texttt{-m}.
  \item Conseguir verificación correcta en modo~0 antes de cualquier otra
        modificación.
  \item Separar los kernels numéricos del código de orquestación (temporizadores,
        E/S, verificación).
  \item Compilar solo los kernels y pasar tamaños y constantes como argumentos.
  \item Cambiar el indexado NumPy encadenado (\texttt{u[k][j][i]}) por indexado
        directo (\texttt{u[k, j, i]}).
  \item Añadir anotaciones de tipo Cython y memoryviews en el modo~3.
\end{enumerate}

La verificación oficial del benchmark fue el criterio de aceptación. Una
optimización solo se mantenía si seguía dando \texttt{SUCCESSFUL}, y las
versiones más rápidas que no verificaban se descartaban. Las decisiones de
implementación que este criterio exigió se detallan en la
Sección~\ref{sec:decisiones-implementacion}.


\section{Decisiones de implementación}
\label{sec:decisiones-implementacion}

\subsection{Separación de kernels y orquestación}
Las funciones de alto nivel de los NAS mezclan inicialización, temporizadores,
E/S, selección de clase, verificación y kernels numéricos. Compilar esas
funciones completas con OMP4Py resulta frágil porque la transformación AST no
puede aplicarse de forma segura a código que combina efectos secundarios con
regiones paralelas. Por ello, la orquestación se mantuvo en Python puro y solo
se compilaron las funciones con carga numérica significativa.

\subsection{Compilación dinámica tras inicialización}
En varios benchmarks, las constantes y tamaños de array dependen de la clase
seleccionada y se calculan durante la ejecución. Si un kernel compilado captura
esos valores antes de que la inicialización se complete, puede usar dimensiones
incorrectas. Para evitarlo, la decoración OMP4Py se aplica después de completar
la inicialización de constantes, en una función \texttt{compile\_kernels()} que
se invoca explícitamente una vez conocida la clase.

\subsection{Indexado directo en NumPy}
Las expresiones como \texttt{u[k][j][i][m]} generan objetos intermedios en
NumPy y dificultan la generación de código eficiente por parte de Cython. La
sustitución por \texttt{u[k, j, i, m]} es semánticamente equivalente y permite
a Cython generar accesos directos a los buffers C cuando se combina con
memoryviews.

\subsection{Variables privadas explícitas}
En EP apareció un fallo de verificación en modo compilado tipado. Varias
variables temporales se compartían entre hilos y producían resultados
incorrectos. La corrección consistió en declararlas como privadas dentro de la
región paralela. Este caso ilustra que una traducción estructuralmente parecida
al código C/OpenMP no garantiza una semántica equivalente en Python. La
visibilidad de las variables locales depende del ámbito léxico del intérprete,
no de las reglas de ámbito de OpenMP.

\subsection{Funciones auxiliares compilables localmente}
Antes de la campaña final se corrigieron varios cuellos de botella derivados de
funciones auxiliares definidas fuera del propio fichero del benchmark. En EP, la
función \texttt{vranlc}, encargada de generar números pseudoaleatorios, se
llamaba como función Python externa desde el kernel compilado, lo que impedía
que el modo~2 compilase todo el cuerpo relevante. La solución fue incluir en
\texttt{EP\_Python.py} versiones localmente compilables de \texttt{vranlc}, una
sin tipos y otra tipada. En IS se revisó la inicialización de la secuencia de
claves para que cada modo ejecutase la ruta correcta y no quedase una parte
secuencial accidental en modo híbrido. Junto con la eliminación del estado
global en MG que se describe más adelante, estas correcciones separan los
límites del port de los límites del runtime. En las versiones finales, los
kernels críticos están dentro del propio fichero del benchmark y son compilables
por OMP4Py/Cython.


\section{Particularidades por benchmark}
\label{sec:particularidades-benchmarks}

\subsection{EP}
EP genera pares de variables aleatorias con distribución normal mediante el
método de Box-Muller. El kernel paralelo central recorre $2^{nk}$ bloques de
dos muestras cada uno. Dentro de cada bloque calcula una semilla local, evalúa
raíces cuadradas y aritmética de punto flotante en escalares, y acumula dos
contadores globales (\texttt{sx}, \texttt{sy}) con una reducción.

La implementación tiene dos versiones compiladas, \texttt{ep\_compute} (modo~2)
y \texttt{ep\_compute\_types} (modo~3). Ambas usan la misma directiva
\texttt{parallel reduction(+:sx,sy)} con planificación estática, pero en
\texttt{ep\_compute\_types} todas las variables del bucle interno están
declaradas con tipos C de Cython (\texttt{cython.int}, \texttt{cython.double})
y se declaran explícitamente privadas para evitar condiciones de carrera. El
módulo principal selecciona la versión tipada en modo~3 mediante
\texttt{use\_compiled\_types()}. Sin ese tipado, cada operación aritmética pasa
por el mecanismo de despacho de objetos Python. La magnitud de esta diferencia
se cuantifica en el Capítulo~\ref{chap:resultados}.

La sección crítica al final de la región paralela actualiza el histograma de
pares. Esta operación es secuencial y ocupa un porcentaje despreciable del
tiempo total en modo~3.

\subsection{CG}
CG resuelve un sistema lineal disperso con el método del gradiente conjugado y
calcula el autovalor más pequeño de la matriz. El bucle de gradiente conjugado
tiene 25 iteraciones. Cada una realiza un producto matriz-vector disperso
seguido de varias operaciones de reducción.

La implementación distingue entre modo no compilado y modo compilado. En los
modos~0 y~1, la función \texttt{conj\_grad} usa directivas
\texttt{parallel~for} independientes para cada bucle. En los modos~2 y~3, la
función \texttt{conj\_grad\_compiled} envuelve todo el bucle de iteraciones en
una única directiva \texttt{parallel}, de modo que el grupo de hilos se crea
una sola vez por llamada y los bucles internos se distribuyen con
\texttt{for}. Esta distinción reduce el overhead de creación y destrucción de
hilos, que es especialmente relevante en clase~S donde los tiempos de kernel
son del orden de décimas de segundo.

Los kernels auxiliares (\texttt{normalize\_z\_compiled},
\texttt{shift\_column\_indices\_compiled}, \texttt{fill\_vector\_compiled})
están anotados con memoryviews de Cython (\texttt{cython.long[:]},
\texttt{cython.double[:]}) para el acceso a arrays. La versión tipada evita que
las operaciones escalares de los bucles internos pasen por objetos Python. El
efecto de esta decisión sobre el rendimiento absoluto se cuantifica en el
Capítulo~\ref{chap:resultados}.

Desde el punto de vista del port, el rasgo más importante de CG es que cada
iteración combina trabajo disperso con varias reducciones escalares. El impacto
de esa estructura sobre la escalabilidad se analiza en los
Capítulos~\ref{chap:resultados} y~\ref{chap:profiling}.

\subsection{IS}
IS ordena enteros de 32 bits mediante un bucket sort en dos fases. La primera
construye el histograma de claves y la segunda realiza la ordenación local en
cada cubo. Para inicializar el vector de claves, el módulo principal usa
\texttt{create\_seq} (función Python sin compilar, modos~0/1/2) o
\texttt{create\_seq\_types} (compilada con Cython y tipos explícitos, modo~3).
Además, las rutas críticas de ordenación usan variables privadas y tipos Cython
explícitos en modo~3. El efecto de esta distinción se cuantifica en el
Capítulo~\ref{chap:resultados}.

La función \texttt{rank} usa una única directiva \texttt{parallel} con un
bloque de variables privadas explícitas por hilo. Cada hilo gestiona su
identificador, los offsets de cubos, las posiciones de inserción y los
contadores locales. Una directiva
\texttt{barrier} explícita separa la fase de histograma de la fase de
ordenación. Todos los hilos deben completar el conteo global de claves antes de
calcular los offsets acumulados, que se calculan en una sección secuencial.

El paso de prefijos acumulados queda como una fase secuencial entre las dos fases
paralelas. En este capítulo interesa como restricción de implementación, y su peso
en tiempo y escalabilidad se estudia después.

\subsection{MG}
MG aplica un ciclo en V del método multigrid para resolver una ecuación de
Poisson tridimensional discreta. Los kernels del ciclo son \texttt{psinv}
(suavizado), \texttt{resid} (cálculo de residuo), \texttt{rprj3} (restricción),
\texttt{interp} (interpolación), \texttt{comm3} (comunicación de fronteras) y
\texttt{zero3} (inicialización).

Durante el port se eliminaron dependencias de estado global en kernels como
\texttt{zran3} y \texttt{norm2u3}, y se añadieron rutas tipadas para que Cython
pudiese compilar los accesos a malla y los acumuladores escalares sin atravesar
el intérprete. La diferencia entre modo~2 y modo~3 es muy grande en tiempo
absoluto, como recoge el Capítulo~\ref{chap:resultados}.

Los arrays de malla se representan como vectores 1D con indexado manual
(\texttt{r[(i3*n2 + i2)*n1 + i1]}) porque las memoryviews de Cython en sus
versiones más accesibles solo funcionan eficientemente con arrays de forma fija
o con vistas 1D. El bucle más externo de cada kernel (\texttt{for i3 in range})
se distribuye entre hilos con \texttt{with omp("for")}.

La particularidad de MG es la jerarquía de niveles. El ciclo en V desciende hasta
una malla de $2 \times 2 \times 2$. En los niveles gruesos, el bucle externo
tiene solo unas pocas iteraciones (1 o 2), lo que hace imposible distribuir
trabajo entre muchos hilos. El efecto de esa granularidad desigual se cuantifica
en el Capítulo~\ref{chap:profiling}.

\subsection{FT}
FT calcula la transformada de Fourier tridimensional discreta de un array
complejo mediante FFTs 1D sucesivas sobre cada eje (tres pasadas). La función
principal abre una única directiva \texttt{parallel} que agrupa todos los
hilos para el cómputo completo. Dentro de esa región, cada pasada llama a
\texttt{cffts1/2/3}, que distribuye las filas del array entre hilos con
\texttt{for} y delega el cómputo a \texttt{cfftz} y \texttt{fftz2}. El patrón
de región \texttt{parallel} exterior persistente evita crear y destruir el
grupo de hilos en cada pasada de la FFT.

La implementación tiene una jerarquía completa de funciones tipadas,
\texttt{fft\_types}, \texttt{cffts1/2/3\_types}, \texttt{cfftz\_types},
\texttt{fftz2\_types}. En sus versiones no tipadas, las operaciones sobre pares
complejos usan el tipo \texttt{complex} de Python, que almacena parte real e
imaginaria como objetos separados y llama a funciones Python para la suma y el
producto. En las versiones tipadas, esos valores son \texttt{cython.doublecomplex}
(\texttt{double complex} de C99), y las operaciones aritméticas son instrucciones
de punto flotante del procesador. Esta diferencia es especialmente visible
porque cada iteración del bucle interior de \texttt{fftz2} realiza dos
multiplicaciones y dos sumas sobre complejos. El impacto se reporta en los
capítulos de resultados.

La pasada sobre el eje $x$ debe completarse antes que la del eje $y$, y esta
antes que la del eje $z$. Dentro de cada pasada las filas son independientes y el
\texttt{parallel~for} funciona, pero la separación entre pasadas introduce
barreras naturales que se analizan en el Capítulo~\ref{chap:profiling}.
